% EMATM0067 — Group Report
% Task 3: Comparative Corpus Analysis
% NeurIPS 2026 template

\documentclass{article}

\usepackage[final]{neurips_2026}
\usepackage[utf8]{inputenc}
\usepackage{amsmath}
\usepackage{amssymb}
\usepackage{graphicx}
\usepackage{booktabs}
\usepackage{float}
\usepackage{hyperref}
\usepackage{multicol}

% Reduce whitespace around floats globally
\setlength{\textfloatsep}{6pt plus 2pt minus 2pt}
\setlength{\floatsep}{4pt plus 2pt minus 2pt}
\setlength{\intextsep}{4pt plus 2pt minus 2pt}
\setlength{\abovecaptionskip}{4pt}
\setlength{\belowcaptionskip}{2pt}

% Fix for NeurIPS 2026 undefined \@trackname in final mode
\makeatletter
\providecommand{\@trackname}{}
\makeatother

\title{Comparative Corpus Analysis of Scientific Abstracts\\
       in AI, Machine Learning and NLP}

\author{
  Ayush Bhadauria \\
  Student ID: 2685201 \\
  \texttt{zj25233@bristol.ac.uk}
  \And
  Jingting Gu \\
  Student ID: 2793772 \\
  \texttt{yn25187@bristol.ac.uk}
  \And
  Pavithraa Gangadhar \\
  Student ID: 2757887 \\
  \texttt{nc25293@bristol.ac.uk}
  \AND
  Chenyang Yan \\
  Student ID: 2765299 \\
  \texttt{th25186@bristol.ac.uk}
  \And
  Balarammahanti Siva Ganesh \\
  Student ID: 2799846 \\
  \texttt{en25786@bristol.ac.uk}
  \AND
  University of Bristol \quad EMATM0067 --- Spring 2026
}

\begin{document}

\maketitle

% ============================================================
% 1. ABSTRACT
% ============================================================
\begin{abstract}
This project explores the development of scientific
research on AI, ML, and NLP through the analysis of
a large corpus of scientific abstracts published from
1991 to 2021. The main objective is to study how the
focus of research, language used, and means of communication
have changed over various eras, and how different text
analytics approaches can influence insights obtained.

In this work, we built a pipeline for text analytics
which consists of such elements as preprocessing, temporal
segmentation, three types of text representation (TF-IDF
unigrams, n-gram TF-IDF, and BERT embeddings), and analytical
approaches such as keyword analysis, LDA, BERTopic, and
dimensionality reduction using Truncated SVD and t-SNE.
Our dataset was divided into four time periods (1991-1998,
1999-2006, 2007-2014, and 2015-2021) in order to find
patterns across \texttt{cs.AI}, \texttt{cs.LG},
and \texttt{cs.CL} subfields.

The results indicate changes in language and emphasis
over time in research, with increasing specialization and
the appearance of sophisticated multiple-word technical
terms in the more recent eras. TF-IDF and n-gram keywords
proved most accurate in pinpointing paradigm shifts. BERTopic
found that the current period is characterized by
theme diversity rather than singular focus on a topic.
LDA produced coherent latent themes but was affected by
vocabulary shift. BERT embeddings arranged documents into
themes rather than eras, indicating a valuable negative
finding for historical research. The results emphasize the
need for balance between interpretability and semantic
depth in methodology, and prove that utilizing multiple
approaches provides a fuller understanding of research evolution.
\end{abstract}


% ============================================================
% 2. INTRODUCTION
% ============================================================
\section{Introduction}

Rapid expansion within the domains of Artificial Intelligence,
Machine Learning, and Natural Language Processing has
resulted in drastic changes within the field of science
communication over time. Abstracts are brief descriptions of
scientific findings that make for an ideal resource to study
this evolution. The analysis of abstracts from different eras
would help identify trends in research focus, terminology,
and methods.
\citep{arxiv2020}.

Understanding the importance of these changes is very clear for many reasons. To begin with, it helps researchers understand the changes that occur in paradigms, for example, from rules-based approaches to data-based approaches and deep learning.
\citep{devlin-etal-2019-bert} — and gives insights into the development of the intricacies and stylistics of scientific discourse as the field progresses.
The study analyzes an extensive number of scientific abstracts derived from the arXiv dataset found in HuggingFace.
\citep{arxiv2020}, focusing specifically on three subject
categories: \texttt{cs.AI} (Artificial Intelligence),
\texttt{cs.LG} (Machine Learning), and \texttt{cs.CL}
(Computing  Language/NLP). The dataset has metadata containing information on the
titles, category, and publication ID, where information
regarding time period was extracted. The corpus spans four
time periods ranging from 1991 to 2021, providing an opportunity
for conducting research longitudinally based on the following
time periods: 1991-1998 (symbolic AI era), 1999-2006
(statistical learning era), 2007-2014 (deep learning)

A major challenge with the problem comes from the aspect
that language undergoes evolution as well. There could
be shifts in terminologies, introduction of new theories,
and sometimes, even when talking about the same
concept, the phrase could differ from one era to another.
Another key issue is the imbalanced distribution
of the data across time spans, where modern years have
many more publications than old years.
\citep{scikit-learn, nltk}.

For instance, consider how the language used for describing qualitative observations differs.
The abstract from the early 1990s could use terms like “expert systems,” “knowledge bases,” and “logic programming”; however, the abstract from 2015 to 2021 would refer to pre-trained transformers.
\citep{devlin-etal-2019-bert}, attention mechanisms,
or convolutional neural networks.

To address the problems stated above,
the process of text analysis becomes systematic and
different techniques for analysis are examined. This time,
however, the task is not only to find any patterns in the
text but also realize how the selection of the technique affects
the information acquired. The two experimental parameters include \textbf{text representation}: we compare
TF-IDF unigrams \citep{scikit-learn}, n-gram TF-IDF
(bigrams and trigrams), and contextual BERT embeddings
\citep{bge2023}. The second is \textbf{topic discovery
method}: we compare TF-IDF keyword analysis, Latent
Dirichlet Allocation (LDA) \citep{blei2003latent}, and
BERTopic \citep{grootendorst2022bertopic}. We additionally
apply Truncated SVD \citep{scikit-learn} to trace the
vocabulary trajectory of the field and t-SNE
\citep{vandermaaten2008} to visualise how BERT embeddings
distribute across periods.
Section~\ref{sec:methods} describes the full pipeline.
Section~\ref{sec:evaluation} presents results and
evaluation. We conclude in Section~\ref{sec:conclusions}.


% ============================================================
% 3. METHODS
% ============================================================
\section{Methods}
\label{sec:methods}

\subsection{Data Preparation}

The corpus was loaded from the arXiv HuggingFace dataset
\citep{arxiv2020} and filtered to \texttt{cs.AI},
\texttt{cs.LG}, and \texttt{cs.CL}.Excess columns were stripped away, papers with incomplete
abstracts were omitted, and category names were made uniform.
Next, papers were grouped into one of four epochs, as
follows: 1991-1998, 1999-2006, 2007-2014, and 2015-2021.
A clear class imbalance is evident in all epochs. The
number of papers in Period 4 greatly exceeds that in other
epochs, due to the exponential growth seen since 2012.
To ensure fairness, keyword and topic frequencies were
standardised rather than taken at face value.

\subsection{Text Preprocessing}

All abstracts were processed as follows:
lowercase conversion, elimination of non-alphabetic
characters and numbers, tokenization, filtering out of
stopwords using the scikit-learn and NLTK English stopwords lists
\citep{scikit-learn, nltk}, followed by
lemmatization based on WordNet with POS tagging for correct
identification of the base forms of verbs, nouns and adjectives.
POS tags were converted to WordNet tags before
lemmatization, which resulted in better accuracy than
the tag-unaware strategy. A custom domain stopword list was additionally
applied to suppress high-frequency but uninformative
academic terms such as ``paper'', ``proposed'', and
``results'', which would otherwise dominate TF-IDF scores
without providing era-discriminative signal.

\subsection{Analytic Axis 1: Text Representation}
\label{subsec:representation}

We compare three text representations, each offering a
different trade-off between interpretability and semantic
richness.

\paragraph{TF-IDF Unigrams (Baseline).}
TF-IDF weights each word by its frequency in a document
relative to the corpus, producing a sparse vector where
each dimension maps directly to a human-readable word.
\textit{Hypothesis: Unigram TF-IDF will identify
distinctive keywords per period but will miss multi-word
methodological phrases and semantically related terms.}

\paragraph{N-gram TF-IDF.}
We extend TF-IDF to bigrams and trigrams (range 2--3)
using \texttt{CountVectorizer} \citep{scikit-learn},
capturing compound phrases like ``convolutional neural
network'' that unigram analysis splits. N-gram density
is also tracked as a proxy for the technical specificity
of writing over time.
\textit{Hypothesis: N-gram TF-IDF will reveal
method-specific vocabulary shifts more clearly than
unigrams, at the cost of a larger, sparser feature space.}

\paragraph{BERT Embeddings.}
Two pre-trained embedding models were used for different
purposes. For semantic similarity analysis and t-SNE
visualisation, the \texttt{all-MiniLM-L6-v2} model
\citep{reimers-2019-sentence-bert} was used to produce
384-dimensional embeddings, chosen for its efficiency on
large corpora. For BERTopic topic modelling, the
\texttt{BAAI/bge-base-en-v1.5} model \citep{bge2023}
was used, selected for its stronger semantic capture
on academic text. Both models compute embeddings as the
mean-pooled final hidden layer over all tokens.
\textit{Hypothesis: BERT embeddings will capture semantic
similarity effectively but may be less sensitive to the
vocabulary-level paradigm shifts that characterise era
transitions.}

\subsection{Analytic Axis 2: Topic Discovery Method}
\label{subsec:comparison}

\paragraph{TF-IDF Keyword Analysis (Baseline).}
For each period and category we extract the top-$k$ terms
by normalised TF-IDF score at both unigram and n-gram
levels, to identify concepts gaining or losing prominence.
\textit{Hypothesis: Keyword analysis will be the most
interpretable approach but may miss latent themes spanning
multiple vocabulary items.}

\paragraph{Period Classification.}
A Logistic Regression classifier \citep{scikit-learn} was
trained to predict the time period label of an abstract
from its TF-IDF representation (max 5,000 features), using
an 80/20 train/test split with \texttt{random\_state=42}.
If there has been much evolution in language,
then it follows that the classifier will perform well. Confusion between adjacent periods suggests similarity in language.
\textit{Hypothesis: Classification accuracy is expected to be higher in discriminating Period 1 from Period 4, and lower between adjacent periods.}

\paragraph{Latent Dirichlet Allocation (LDA).}
LDA \citep{blei2003latent} The model represents each document as a mixture of hidden topics, and each topic as a distribution over words.
We use this approach on our preprocessed data set to identify the themes which reoccur in our data without using any labeling information.
\textit{Hypothesis: LDA will complement keyword analysis by grouping co-occurring terms into coherent themes, though its bag-of-words assumption may limit sensitivity to rapidly evolving technical vocabulary.}

\paragraph{BERTopic.}
BERTopic \citep{grootendorst2022bertopic} clusters documents
by their BERT embeddings and labels each cluster using
class-based TF-IDF (c-TF-IDF). We customise the pipeline
using UMAP \citep{mcinnes2018umap} for dimensionality
reduction based on semantic relationships, and HDBSCAN
\citep{campello2013hdbscan} for density-based clustering
using Euclidean distance, chosen to produce robust,
non-fragmented topic clusters. Unigrams and bigrams are
retained for topic word extraction.Topic Intensity: 
The frequency of occurrence of a particular topic in a period and
category is measured through time to uncover the development of 
research themes within specific subfields.
\textit{Research Hypothesis:
BERTopic will generate topics with greater semantic coherence than 
LDA since BERTopic relies on embeddings and not bag of words 
methodology.}

\subsection{Temporal Drift Visualisation}

\textbf{Truncated SVD} \citep{scikit-learn} reduces the 2D representation of the TF-IDF matrix; centroid points of periods are linked to visualize the trajectory of vocabulary used in the field at the macro level, based on
labels for axes that correspond to the most extreme terms at each end of the axis. \textbf{t-SNE} \citep{vandermaaten2008} This is done by applying
BERT to a random sample of 5,000 abstracts (perplexity 45).


% ============================================================
% 4. EXPERIMENTAL EVALUATION
% ============================================================
\section{Experimental Evaluation}
\label{sec:evaluation}

\subsection{Metrics and Setup}

We evaluate each approach against three criteria:
\textbf{consistency} (do findings replicate across
methods?), \textbf{interpretability} (are outputs
understandable to a domain reader?), and
\textbf{alignment} (do findings match known milestones
such as the rise of deep learning post-2012?).
Interpretability and alignment are assessed qualitatively;
we mitigate subjectivity by cross-referencing findings
across all methods. Table~\ref{tab:summary} summarises
the relative strengths of each approach.

All preprocessing artefacts were saved and reused across
experiments to guarantee consistent inputs. TF-IDF and
n-gram features used \texttt{scikit-learn}
\citep{scikit-learn}. BERT embeddings used the
\texttt{BAAI/bge-base-en-v1.5} model from HuggingFace
\citep{bge2023}. BERTopic used the \texttt{bertopic}
library \citep{grootendorst2022bertopic} with UMAP
\citep{mcinnes2018umap} and HDBSCAN
\citep{campello2013hdbscan} for clustering.

\begin{table}[H]
\centering
\caption{Qualitative comparison of all methods against
evaluation criteria. \checkmark = strong,
$\sim$ = moderate, $\times$ = weak.}
\label{tab:summary}
\begin{tabular}{lccc}
\toprule
\textbf{Method} & \textbf{Consistency} &
\textbf{Interpretability} & \textbf{Alignment} \\
\midrule
TF-IDF Unigrams     & \checkmark & \checkmark & \checkmark \\
N-gram TF-IDF       & \checkmark & \checkmark & \checkmark \\
BERT Embeddings     & $\times$   & $\times$   & $\times$   \\
LDA                 & $\sim$     & $\sim$     & $\sim$     \\
BERTopic            & \checkmark & $\sim$     & \checkmark \\
SVD Trajectory      & \checkmark & $\sim$     & \checkmark \\
\bottomrule
\end{tabular}
\end{table}

\subsection{Corpus Growth and Data Imbalance}

\begin{figure}[H]
  \centering
  \includegraphics[width=0.80\linewidth, trim=0 30pt 0 30pt, clip]{figures/01_category_growth.png}
  \caption{ArXiv AI subfield publication growth (1991--2021).
           cs.LG dominates from 2015, reaching $\sim$37,000
           papers annually by 2021.}
  \label{fig:growth}
\end{figure}

Figure~\ref{fig:growth}indicates nearly flat publication production
till about 2010, after which there is exponential increase
due to cs.LG. This validates the extreme imbalance between
periods: Period 4 is much larger compared to previous periods
by a large margin. As recent papers will completely skew the numbers,
the keyword and topic weights have been normalized everywhere.

\subsection{Axis 1 Results: N-gram Key Terms per Period}

\begin{figure}[H]
  \centering
  \includegraphics[width=\linewidth, trim=0 40pt 0 40pt, clip]{figures/09_multi_phrases.png}
  \caption{Top 10 bigram and trigram key terms per period
           from n-gram TF-IDF analysis.}
  \label{fig:phrases}
\end{figure}

Figure~\ref{fig:phrases} shows the most obvious picture of
how research language developed. Between 1991–1998, common collocations
are \textit{machine translation}, \textit{knowledge base},
\textit{decision tree}, and \textit{speech recognition} —
the vocabulary of symbolic AI and early NLP. By 1999--2006,
\textit{logic program} indicating the combination of symbolic and probabilistic methodologies.
In 2007--2014, \textit{machine learn} takes over alongside
\textit{support vector} and \textit{reinforcement learn},
signaling the shift towards statistics. In 2015–2021,
\textit{machine learn} is overwhelmingly frequent
(count $>15{,}000$), joined by \textit{deep neural},
\textit{convolutional neural}, \textit{pre train}, and
\textit{reinforcement learn} — reflects the language of the deep
learning period. It verifies our hypothesis: n-grams uncover
compound methodological terminology inaccessible to unigrams.

\subsection{Axis 1 Results: Keyword Evolution Heatmap}

\begin{figure}[H]
  \centering
  \includegraphics[width=0.72\linewidth, trim=0 30pt 0 30pt, clip]{figures/08_keywords_period_evolution.png}
  \caption{Rise and decay of AI concepts (normalised TF-IDF).
           Red = high score; purple = absent.}
  \label{fig:evolution}
\end{figure}

Figure~\ref{fig:evolution} tracks 30 terms across all
periods. Symbolic AI terms — \textit{grammar}, \textit{parse},
\textit{logic} — score modestly early and disappear by
Period 4. \textit{Deep} scores 0.0 in the first two periods,
rises to 52.4 in 2007--2014, and reaches 2,114.4 by
2015--2021. \textit{Network} grows from 7.5 to 3,635.9
and \textit{learn} from 19.6 to 4,080.8. The transition
accelerates sharply after 2014, consistent with the
n-gram results. Agreement between both TF-IDF methods
strengthens confidence that these are genuine field-level
shifts rather than noise.

\subsection{Axis 1 Results: Period Classification}

A Logistic Regression classifier was trained to predict
time period from TF-IDF representations. The classifier
achieved an accuracy of \textbf{90.9}\%
on the held-out test set.

This result [confirms / partially confirms] our hypothesis
that temporal linguistic drift is detectable from text
representations alone. The fact that a simple linear
classifier can distinguish periods with meaningful accuracy
provides quantitative evidence that vocabulary-level
changes across the four eras are systematic rather than
random. This finding complements the qualitative evidence
from the keyword analysis and n-gram results.

\subsection{Axis 2 Results: LDA Topic Analysis}

We applied LDA with $k=5$ topics to the preprocessed corpus
using a bigram-aware \texttt{CountVectorizer} with a vocabulary
of 1,000 features. Table~\ref{tab:lda} shows the top keywords
for each discovered topic and our interpretation of each.

\begin{table}[H]
\centering
\small
\caption{LDA topic definitions (top keywords) and
         interpretations.}
\label{tab:lda}
\begin{tabular}{cll}
\toprule
\textbf{Topic} & \textbf{Top Keywords} & \textbf{Interpretation} \\
\midrule
1 & learn, machine, machine learn, & General ML \\
  & optimization, technique, cluster & Methods \\
2 & function, distribution, sample, & Statistical / \\
  & loss, bound, gradient, estimate & Learning Theory \\
3 & agent, policy, human, environment, & Reinforcement \\
  & user, question, reinforcement & Learning / Agents \\
4 & graph, label, representation, train, & NLP / Knowledge \\
  & domain, knowledge, speech, sentence & Representation \\
5 & network, neural, neural network, & Deep Learning / \\
  & image, deep, architecture, layer & Neural Networks \\
\bottomrule
\end{tabular}
\end{table}

\begin{figure}[H]
  \centering
  \includegraphics[width=0.90\linewidth, trim=0 40pt 0 40pt, clip]{figures/lda_category_period_evolution.png}
  \caption{LDA topic intensity across categories and periods.
           Topic 4 (orange) grows in cs.CL; Topic 5 (green)
           rises across all subfields from 2007.}
  \label{fig:lda}
\end{figure}

Figure~\ref{fig:lda} and the underlying data reveal several
clear patterns. In \texttt{cs.CL},Topic 4 (NLP and
knowledge representation) is the prevalent topic that
dominates the entire corpus and steadily increases in all four
time periods, starting from 0.41 in 1991-1998 to 0.58 in
2015-2021. It is the most prominent trend in the LDA output,
indicating the continuous interest in representation learning
and language understanding among the NLP community.
Meanwhile, Topic 1 (general ML techniques) steadily
declines from 0.27 to 0.11 in cs.CL during the same time
period due to the specialisation of the subfield.

In \texttt{cs.AI}, Topic 3 (reinforcement learning and
agent-based models) continues to be a dominant theme
across all periods, having a maximum value of 0.39 between
1999–2006 and staying relatively constant at around 0.36 by
2015–2021. This demonstrates the continued importance
of the agent-based paradigm within the AI subfield, despite
its increasing dominance in other parts of the discipline. Topic
5 (deep learning and neural networks) goes from 0.04
between 1991–1998 to

In \texttt{cs.LG}, Concerning Topic 2 (statistical and mathematical learning theory), there is a pronounced jump upward, reaching from 0.11 in 1991--1998 to 0.31 in 1999--2006 and staying high until 2007--2014. It corresponds to the rise of statistics and the kernel machine learning framework in this time. In 2015--2021, Topic 5 (neural networks) reaches 0.21, corresponding to the migration of cs.LG to deep learning.

The results partially support our hypothesis. LDA was able to detect meaningful thematic clusters corresponding to the evolution of each subfield. Nevertheless, the presence of the common term "learn" in all topics, as a result of its survival after the stopword removal procedure, blurs the discrimination power of the topic clusters. Indeed, it is well-known that the use of LDA on homogeneous-domain corpora is hindered by the saturation of topic-word probability distributions with the central terms of the field, which is why the separation between the topics suffers in comparison with what BERTopic can achieve via embedding clustering.

\subsection{Axis 2 Results: BERTopic Evolution}

\begin{figure}[H]
  \centering
  \includegraphics[width=0.90\linewidth, trim=0 30pt 0 30pt, clip]{figures/bert_category_period_evolution.png}
  \caption{BERTopic topic intensity across categories and
           periods. Topic 1 (teal) peaks in 2007--2014 then
           declines sharply by 2015--2021.}
  \label{fig:bert}
\end{figure}

Figure~\ref{fig:bert} Topic 1 is seen to dominate the beginning and middle stages,
reaching its maximum value of 0.31 in cs.CL (2007-2014)
and 0.30 in cs.LG (1999-2006), after which it declines sharply
below 0.05 during 2015-2021. This is not an indication of a
problem with BERTopic but a sign that there was actual diversification
in the field: As the number of papers grew exponentially,
researchers branched out into computer vision, reinforcement learning,
language modelling.

\subsection{Temporal Drift: SVD and t-SNE}

\begin{figure}[H]
  \centering
  \begin{minipage}[t]{0.48\linewidth}
    \centering
    \includegraphics[width=\linewidth, trim=0 30pt 0 30pt, clip]{figures/04_tfidf_trajectory.png}
    \caption{Truncated SVD projection. Period centroids
             trace a rightward ``buyer $\to$ learn''
             semantic trajectory.}
    \label{fig:trajectory}
  \end{minipage}
  \hfill
  \begin{minipage}[t]{0.48\linewidth}
    \centering
    \includegraphics[width=\linewidth, trim=0 30pt 0 30pt, clip]{figures/03_embeddings_clustering.png}
    \caption{t-SNE of BERT embeddings (5,000 abstracts).
             Periods are interspersed: BERT organises
             by topic, not era.}
    \label{fig:tsne}
  \end{minipage}
\end{figure}

Figure~\ref{fig:trajectory} centroids for periods shift
in the direction to the right, with the biggest jump from
Period 3 to Period 4, supporting the trend of vocabulary
change acceleration since 2014. Figure~\ref{fig:tsne}
reveals that BERT embeddings do not group based on the
era of origin — abstracts from each era are scattered
throughout the 2D space. This is an important negative
result: BERT captures the topic level but not the era
level of research, rendering TF-IDF more susceptible to
vocabulary drift over time.

\subsection{Error Analysis and Discussion}

A judgment had to be made regarding the stop word list because although
words like ``learn'' and ``network'' have discriminating power, their extremely
high frequency in Period 4 could obscure other themes. This had an immediate
impact on LDA because the term ``learn'' appears in all five topic-word
frequency distributions, making the boundaries between the topics
($>1{,}000$ features) Topic separation would be more effective.
The intensity collapse of BERTopic between 2015-2021 indicates
diversification but is incapable of characterizing the actual
content of that diversity without investigating topic label
lists. The t-SNE plot is influenced by class imbalance since
Period 4 is prevalent in the 5,000-abstacts dataset, causing
other time periods to become underrepresented geometrically.

Nevertheless, results have shown high consistency across methods.
TF-IDF keyword and n-gram extraction generated the most
interpretable and reliable results. BERTopic brought an additional
perspective on thematic diversification. The LDA model provided
interpretable topics but was not as consistent with hyperparameters.
BERT embeddings were found to be the least informative in era-level
temporal analysis. One limitation of the approach is that the
categorization method treats each category as static and
non-overlapping while papers often overlap with subfields.


% ============================================================
% 5. TEAM CONTRIBUTIONS
% ============================================================
\section{Team Contributions}
\label{sec:contributions}

\textbf{Pavithraa Gangadhar} developed and established
the complete data pipeline process, including the
processes of loading, evaluating quality, cleaning, and
segmentation of data based on time period. She created the
text pre-processing strategy, which includes tokenisation,
stopwords elimination, and lemmatisation, and developed the
TF-IDF, n-gram, and transformer approaches for the creation
of text representations. Finally, she performed exploratory
data analysis, keyword extraction, similarity analysis, and
classification of data by time period.

\textbf{Jingting Gu} Used BERTopic modelling, applying
the \texttt{BAAI/bge-base-en-v1.5} embedding model,
with UMAP for dimensional reduction and HDBSCAN for density-based topic clustering.
She included POS tagging within the preprocessing stage, employing WordNet tags to enhance
lemmatisation efficiency, and designed a universal interface for all downstream visualization tools.
She also constructed the pipeline orchestration layer that allows for automatic execution of all
preprocessing, modelling, and visualization modules through one single access point.

\textbf{Chenyang Yan} was responsible for designing and implementing the entire
visualization tool suite, incorporating t-SNE for semantic clustering and Truncated SVD for macroscopic trajectory analysis. He devised and optimized the LDA model pipeline himself, crafting a novel context-based keyword extraction system capable of identifying hidden transitions in topics over different time periods. He also created a proprietary stopword list specific to the domain under study.

\textbf{Balarammahanti Siva Ganesh} The data preprocessing pipeline and dataset engineering
process were applied to obtain a clean structured
dataset after dealing with missing data, inconsistencies,
and duplications in the raw dataset. He built
preprocessing scripts that can be reused in the
\texttt{src} directory, created structured features for
analysis, and applied temporal segmentation to set the
baseline distribution of the dataset.

\textbf{Ayush Bhadauria}coordinating the entire project
workflow by distributing tasks among the team members
depending on their strong points and making sure
that everything was included into one workflow.
The conversations with the other members were made in
order to know how they solved the task and what they
obtained; all this was used for writing this report.
All the information obtained during our teamwork was
analyzed and written into the report using the same
methods, evaluation and conclusion techniques.


% ============================================================
% 6. CONCLUSIONS
% ============================================================
\section{Conclusions}
\label{sec:conclusions}

This study employed a multi-method text analytics workflow that was
applied to a set of 30-year arXiv abstracts from the AI, ML, and NLP subfields.
In all methods, the results reveal the same consistent pattern, as the corpus moves
from symbolic AI and formalism in the 1990s, then to probability and statistical
learning in the 2000s, then to neural approaches starting from 2007, up to a very
diverse and transformer-led corpus since 2015.

The best methods to identify this shift were those using TF-IDF keyword and n-gram
analysis. Multi-word terms
such as \textit{convolutional neural}, \textit{pre train},
and \textit{reinforcement learn} were important signals that
could not be found using unigrams. BERTopic introduced an interesting point: the
current age can be described through the fragmentation of scientific research into
many specialized areas, seen from the sharp decrease in intensity per topic since
2014. The LDA algorithm identified interpretable thematic groups that largely
correlated with the keyword trends but was somewhat more dependent on the value of
hyperparameters compared to BERTopic.

The embeddings generated by BERT were not grouped according to historical periods,
and this result can be considered quite useful. Semantic representations capture the
meaning of scientific research at the level of topics, while TF-IDF-based vectors
are more suitable for studying vocabulary change over historical periods.

Possible further steps may include solving the problem of class imbalance,
developing quantitative interpretability metrics, and analyzing not only abstracts but
also papers or citation networks.

Code repository: \url{https://github.com/gyy2333/uob-ematm0067-team23}


% ============================================================
% REFERENCES
% ============================================================
\bibliographystyle{unsrt}
\bibliography{references}

\end{document}